\chapter{Doubly Robust Estimation}
\label{chap:doubly-robust}

\section{The augmented estimator}
\label{sec:dr-aipw}

Chapter~\ref{chap:propensity} estimated the ATE from a model for treatment.
The g formula~\eqref{eq:g-formula} suggests instead a model for the outcome,
$\mu_a(x) = \E\left[ Y \mid A = a, X = x \right]$, giving the outcome regression
estimator\index{outcome regression} $n^{-1}\sum_i \left( \hat{\mu}_1(X_i) - \hat{\mu}_0(X_i) \right)$.
Both estimators require their model to be correct. The augmented estimator
requires only that one of them be.

\begin{definition}[Augmented IPW estimator]
\label{def:aipw}
\index{augmented inverse probability weighting}
\begin{equation}
  \hat{\ATE}_{\mathrm{AIPW}} = \frac{1}{n} \sum_{i=1}^{n} \left\{
    \hat{\mu}_1(X_i) - \hat{\mu}_0(X_i)
    + \frac{A_i \left( Y_i - \hat{\mu}_1(X_i) \right)}{\hat{e}(X_i)}
    - \frac{(1 - A_i)\left( Y_i - \hat{\mu}_0(X_i) \right)}{1 - \hat{e}(X_i)}
  \right\}.
  \label{eq:aipw}
\end{equation}
\end{definition}

\begin{theorem}[Double robustness]
\label{thm:double-robustness}
Let $e^{*}$ and $\mu_a^{*}$ denote the probability limits of $\hat{e}$ and
$\hat{\mu}_a$. If either $e^{*} = e$ or $\mu_a^{*} = \mu_a$ for $a \in \{0,1\}$,
and positivity holds, then $\hat{\ATE}_{\mathrm{AIPW}}$ is consistent for the
ATE.
\end{theorem}

\begin{proof}
Write the population limit of the $a = 1$ portion of~\eqref{eq:aipw} as
\begin{equation*}
  \E\left[ \mu_1^{*}(X) + \frac{A\left( Y - \mu_1^{*}(X) \right)}{e^{*}(X)} \right]
  = \E\left[ Y(1) \right]
    + \E\left[ \left( 1 - \frac{e(X)}{e^{*}(X)} \right)
               \left( \mu_1^{*}(X) - \mu_1(X) \right) \right],
\end{equation*}
which follows by conditioning on $X$, using $\E[A \mid X] = e(X)$ and
$\E[AY \mid X] = e(X)\mu_1(X)$ under Assumptions~\ref{ass:consistency}
and~\ref{ass:ignorability}, and collecting terms. The remainder is a product of
two errors. If $e^{*} = e$ the first factor vanishes. If $\mu_1^{*} = \mu_1$ the
second vanishes. In either case the limit is $\E[Y(1)]$. The argument for
$a = 0$ is identical, and the difference of the two limits is the ATE.
\end{proof}

The proof shows more than the statement. The bias is a product of the two model
errors, so the estimator is not merely robust to one failure but also
second order accurate when both models are somewhat wrong. This is the
observation that makes machine learning nuisance estimation viable.

\begin{corollary}[Rate double robustness]
\label{cor:rate-dr}
If $\|\hat{e} - e\|_{2} = o_p(n^{-r_e})$ and
$\|\hat{\mu}_a - \mu_a\|_{2} = o_p(n^{-r_\mu})$ with $r_e + r_\mu \geq 1/2$,
and cross fitting is used, then $\hat{\ATE}_{\mathrm{AIPW}}$ is
$\sqrt{n}$ consistent and asymptotically normal.
\end{corollary}

\begin{proof}
The remainder term in the proof of Theorem~\ref{thm:double-robustness} is
bounded by $\|\hat{e} - e\|_{2}\,\|\hat{\mu}_a - \mu_a\|_{2}$ after a Cauchy
Schwarz step and the positivity bound on $e^{*}$. Under the stated rates this is
$o_p(n^{-1/2})$ and therefore negligible relative to the influence function
term. Cross fitting removes the empirical process term that would otherwise
require a Donsker condition on the nuisance estimators.
\end{proof}

Corollary~\ref{cor:rate-dr} is proved in greater generality by
\citet{barros2019rates}. It is the reason a random forest with a convergence
rate of $n^{-1/4}$ in each nuisance function suffices for valid inference on the
ATE, whereas the same forest plugged into the g formula alone does not.

\section{Cross fitting}
\label{sec:dr-crossfit}

Cross fitting\index{cross fitting} estimates the nuisance functions on data
disjoint from the fold at which they are evaluated \citep{lindqvist2016crossfit}. Without it, an overfitted
$\hat{\mu}$ correlates with the residual it is subtracted from, and the
correction term in~\eqref{eq:aipw} no longer has mean zero.

\begin{algorithm}[t]
\SetAlgoLined
\KwIn{Data $\{(Y_i, A_i, X_i)\}_{i=1}^{n}$, folds $K$, learners for $e$ and $\mu_a$}
\KwOut{$\hat{\ATE}_{\mathrm{AIPW}}$ and its standard error}
partition the index set into folds $I_1, \dots, I_K$ of near equal size\;
\For{$k \leftarrow 1$ \KwTo $K$}{
  fit $\hat{e}^{(-k)}$ and $\hat{\mu}_a^{(-k)}$ on the units outside $I_k$\;
  \ForEach{$i \in I_k$}{
    evaluate $\hat{e}^{(-k)}(X_i)$ and $\hat{\mu}_a^{(-k)}(X_i)$\;
    truncate $\hat{e}^{(-k)}(X_i)$ into $[\delta, 1 - \delta]$\;
    form the influence contribution $\hat{\psi}_i$ from~\eqref{eq:influence}\;
  }
}
$\hat{\ATE}_{\mathrm{AIPW}} \leftarrow n^{-1} \sum_i \hat{\psi}_i$\;
$\widehat{\mathrm{se}} \leftarrow n^{-1/2}$ times the sample standard deviation
of $\hat{\psi}_i$\;
\caption{Cross fitted augmented estimation of the average treatment effect.}
\label{alg:crossfit}
\end{algorithm}

The truncation step\index{truncation} in Algorithm~\ref{alg:crossfit} is not cosmetic. A learner
that returns $\hat{e} = 0.998$ for a control unit creates a single influence
contribution of order 500 times the others, and the sample standard deviation
that supplies the standard error is then estimating the fourth moment of the
weight distribution rather than anything about the effect. The value
$\delta = 0.01$ is conventional. Reporting the number of truncated units is
obligatory, and if that number is more than a few percent the estimand should be
changed rather than the truncation loosened.

\section{The efficient influence function}
\label{sec:dr-influence}

\begin{proposition}[Efficient influence function]
\label{prop:eif}
\index{influence function}
The efficient influence function for the ATE in the nonparametric model is
\begin{equation}
  \psi(Y, A, X) = \mu_1(X) - \mu_0(X)
    + \frac{A\left( Y - \mu_1(X) \right)}{e(X)}
    - \frac{(1-A)\left( Y - \mu_0(X) \right)}{1 - e(X)}
    - \ATE,
  \label{eq:influence}
\end{equation}
and the semiparametric efficiency bound\index{efficiency bound} is $\Var\left[ \psi \right] / n$.
\end{proposition}

The bound decomposes usefully. Writing $\sigma_a^{2}(x) = \Var\left[ Y \mid A = a, X = x \right]$,
\begin{equation}
  \Var\left[ \psi \right]
    = \E\left[ \frac{\sigma_1^{2}(X)}{e(X)} + \frac{\sigma_0^{2}(X)}{1 - e(X)} \right]
    + \Var\left[ \mu_1(X) - \mu_0(X) \right].
  \label{eq:efficiency-bound}
\end{equation}
The first term is irreducible sampling noise inflated by the imbalance of the
design, and it is what explodes when $e$ approaches zero or one. The second is
the genuine heterogeneity of the effect across covariate strata and is
irreducible in a different sense: it would remain in a randomised experiment.

\begin{table}[t]
\centering
\caption[Estimators compared on the retraining data]{Estimators compared on the
retraining data. The nuisance functions for the last three rows are fitted with
a gradient boosted ensemble under five fold cross fitting. The final column is
the ratio of the estimated variance to the efficiency bound
of~\eqref{eq:efficiency-bound} evaluated at the cross fitted nuisances.}
\label{tab:dr-comparison}
\begin{tabular}{lrrrr}
\toprule
Estimator & Estimate & Std. error & 95\% interval & Var. ratio \\
\midrule
Outcome regression, linear   & 1372 & 174 & (1031, 1713) & 1.41 \\
Hajek weighting              & 1461 & 388 & (700, 2222)  & 6.99 \\
AIPW, both models linear     & 1409 & 191 & (1035, 1783) & 1.69 \\
AIPW, cross fitted ensemble  & 1418 & 147 & (1130, 1706) & 1.01 \\
AIPW, score misspecified     & 1431 & 152 & (1133, 1729) & 1.08 \\
AIPW, outcome misspecified   & 1404 & 163 & (1085, 1723) & 1.24 \\
\bottomrule
\end{tabular}
\end{table}

Table~\ref{tab:dr-comparison} illustrates Theorem~\ref{thm:double-robustness}
directly. The last two rows deliberately break one nuisance model each, by
fitting the score with an intercept only model and by fitting the outcome model
on a single covariate respectively. Both remain within 2 percent of the fully
specified estimate, whereas the pure weighting and pure regression estimators
differ from each other by 6 percent and from the doubly robust estimate in
opposite directions.

The variance ratio column carries the other message of this chapter. The cross
fitted ensemble attains the efficiency bound. The Hajek estimator is seven times
the bound, which means that six sevenths of the sampling variability in that row
is a property of the estimator rather than of the data.
